\chapter{Algorithms and code}
\label{chap:algorithms}

\section{Pseudocode}

\Cref{alg:dijkstra} is the shortest-path algorithm by Dijkstra~\cite{dijkstra1959}.

\begin{algorithm}[htbp]
  \caption{Dijkstra's shortest paths}
  \label{alg:dijkstra}
  \begin{algorithmic}[1]
    \Require graph $G=(V,E)$, weights $w \ge 0$, source $s$
    \Ensure distances $d[v]$ for every $v \in V$
    \State $d[v] \gets \infty$ for all $v$; \quad $d[s] \gets 0$
    \State $Q \gets V$
    \While{$Q \neq \emptyset$}
      \State $u \gets \argmin_{v \in Q} d[v]$; \quad $Q \gets Q \setminus \{u\}$
      \ForAll{$(u,v) \in E$}
        \If{$d[u] + w(u,v) < d[v]$}
          \State $d[v] \gets d[u] + w(u,v)$
        \EndIf
      \EndFor
    \EndWhile
    \State \Return $d$
  \end{algorithmic}
\end{algorithm}

\section{Source code}

\Cref{lst:rust} is a simplified version of how \rutex{} launches a build. Every
build gets its own process group, so nothing is left running when you quit.

\begin{lstlisting}[language=Rust, caption={Spawning latexmk without leaving orphans.}, label={lst:rust}]
pub fn command(program: &str) -> Command {
    let mut cmd = Command::new(program);
    unsafe {
        cmd.pre_exec(|| {
            libc::setpgid(0, 0); // own process group
            libc::prctl(libc::PR_SET_PDEATHSIG, libc::SIGKILL);
            Ok(())
        });
    }
    cmd
}
\end{lstlisting}

\begin{lstlisting}[language=Python, caption={The same idea in Python.}, label={lst:python}]
import os, subprocess

def run(args):
    # start_new_session puts the child in its own process group
    return subprocess.Popen(args, start_new_session=True)
\end{lstlisting}
